\section{Datasets, Simulators and Toolkits}
\label{sec:datasets}

Progress in acoustic perception for autonomous systems is limited less by models than by data recorded on real platforms.
\Cref{tab:datasets} lists the main public resources and what each can test.

\begin{table*}[t]
\centering
\footnotesize
\setlength{\tabcolsep}{3pt}
\begin{tabular}{@{}L{3.1cm}L{2.0cm}L{1.3cm}L{1.4cm}L{7.6cm}@{}}
\toprule
Resource & Platform & Real ego-noise & Ground truth & What it tests \\
\midrule
DREGON~\cite{strauss2018dregon} & drone, 8 mics & yes & motion capture & in-flight localization; per-rotor speed and IMU logged \\
SP Cup 2019~\cite{deleforge2019spcup} & drone & yes & yes & localization for search and rescue \\
SP Cup ego-noise~\cite{deleforge2019spcup} & 10 drones, 1--16 mics & yes & -- & ego-noise of different drones; used in \cref{sec:case} \\
AIRA-UAS~\cite{ruizespitia2018aira} & drone & yes & -- & ego-noise during maneuvers and from nearby drones \\
AVQ~\cite{wang2019avq} & drone, 8 mics + camera & yes & from video & audio-visual localization, tracking and enhancement outdoors \\
LOCATA~\cite{lollmann2018locata,evers2020locata} & robot head (12 mics), spherical, planar, hearing aids; static and moving & no & motion capture & localization and tracking of moving talkers in a reverberant room; used in \cref{sec:case} \\
OVAD~\cite{schulz2021hearing} & vehicle, 56 mics & partly & yes & detection of occluded vehicles \\
STARSS23~\cite{shimada2023starss23} & static array + 360$^\circ$ video & no & motion capture & audio-visual SELD in real rooms \\
AudioSet~\cite{gemmeke2017audioset} & web video & no & labels & event classification and pretraining \\
ESC-50~\cite{piczak2015esc} & curated clips & no & labels & environmental sound classification \\
FSD50K~\cite{fonseca2022fsd50k} & curated clips & no & labels & open event classification \\
VGGSound~\cite{chen2020vggsound} & web video & no & labels & audio-visual recognition \\
\midrule
SoundSpaces~\cite{chen2020soundspaces,chen2022soundspaces2} & simulator & no & exact & audio-visual navigation and rendering \\
Sonicverse~\cite{gao2023sonicverse} & simulator & no & exact & multisensory household agents \\
ObjectFolder~\cite{gao2021objectfolder} & object models & no & exact & visual, auditory and tactile object properties \\
\midrule
HARK~\cite{nakadai2017hark} & toolkit & -- & -- & robot localization, separation, recognition \\
ODAS~\cite{grondin2022odas} & toolkit & -- & -- & embedded localization, tracking, separation \\
\bottomrule
\end{tabular}
\caption{Public datasets, simulators and toolkits. Only DREGON combines real ego-noise with the platform's own motor state and ground-truth directions; none yet combines a moving platform, real ego-noise, platform state, high-rate vibration sensing, multiple sound classes and video.}
\label{tab:datasets}
\end{table*}

Three gaps stand out.
First, the benchmarks with rich annotations (STARSS23 for SELD, SoundSpaces for navigation, AudioSet for classification) do not include ego-noise, while the datasets that do (DREGON, AIRA-UAS, AVQ) are small and focus on localization and enhancement.
Second, no public benchmark yet combines a moving platform, real ego-noise, the platform's state, high-rate vibration sensing, multiple sound classes and visual input.
Such a benchmark would let the embodied-learning and robot-audition communities measure progress on the same problem.
Third, availability is fragile: of the ten drone ego-noise sets released with the SP Cup, two could no longer be downloaded at the time of writing, so archival hosting with persistent identifiers matters as much as collection.
Toolkits such as HARK and ODAS lower the barrier to real-time deployment, but they implement classical pipelines and would benefit from integration with learned models.
